\begin{table}[htbp!]
\centering
\caption{Resultados da auditoria do \textit{dataset}}
\label{tab:auditoria}
\small
\begin{tabular}{|p{7.3cm}|c|c|}
\hline
\textbf{Verificação} & \textbf{Antes da adjudicação} & \textbf{Após} \\
\hline
\multicolumn{3}{|l|}{\textit{Checagens automáticas}} \\ \hline
Problemas estruturais (ids, duplicatas, enumerados, datas) & --- & 0 \\ \hline
Consultas da camada P conferidas na linha de origem & --- & 22/22 \\ \hline
Apontamentos do anotador por regras sem resolução & --- & 0 \\ \hline
Expressões com anotação inconsistente sem justificativa & --- & 0 \\ \hline
\multicolumn{3}{|l|}{\textit{Anotação independente às cegas (310 consultas)}} \\ \hline
Decisão (chamar / não chamar / observacional) & 302/310 (97,4\%) & 310/310 (100,0\%) \\
\quad kappa de Cohen & 0,658 & 1,000 \\ \hline
Leituras compatíveis nos dois sentidos & 300/310 (96,8\%) & 309/310 (99,7\%) \\ \hline
Leitura preferencial idêntica$^{a}$ & 294/298 (98,7\%) & 294/298 (98,7\%) \\ \hline
Presença de cada campo na leitura preferencial$^{a,b}$ (kappa) & 0,999 & 0,999 \\ \hline
Valor do campo, quando ambos o anotam$^{a}$ & 534/537 (99,4\%) & 534/537 (99,4\%) \\ \hline
Consulta com mais de uma leitura aceita$^{a}$ & 296/298 (99,3\%) & 297/298 (99,7\%) \\
\quad kappa de Cohen & 0,981 & 0,990 \\ \hline
Consultas com divergência em alguma medida, sem decisão registrada & 14 & 0 \\ \hline
\end{tabular}
\fonte{Elaborada pelo autor a partir de \texttt{data/auditoria/}. (a) Sobre as consultas em que gabarito e anotador pedem a chamada da ferramenta nas métricas principais. (b) Doze campos por consulta.}
\end{table}
